\setcounterpageref{ex}{0}
\TaoPhieuKiemTra{07}

\setcounter{ex}{0}
\setcounter{bt}{0}

\textbf{\large I. PHẦN TRẮC NGHIỆM (3,0 điểm)}

{\itshape Khoanh tròn vào chữ cái đặt trước câu trả lời đúng. Mỗi câu đúng được 0,5 điểm.}

\begin{ex} % Câu 1
Số bé nhất có hai chữ số khác nhau là
\choice
{\True $10$}
{$11$}
{$12$}
{$01$}
\loigiai{
Số bé nhất có 2 chữ số khác nhau là $10$.\\
\textbf{Đáp án đúng: A}
}
\end{ex}

\begin{ex} % Câu 2
Trong phép tính $87 - 35 = 52$, số $52$ được gọi là
\choice
{Số bị trừ}
{Số trừ}
{\True Hiệu}
{Tổng}
\loigiai{
Kết quả của phép trừ được gọi là Hiệu.\\
\textbf{Đáp án đúng: C}
}
\end{ex}

\begin{ex} % Câu 3
Tính: $8\text{ }l + 6\text{ }l = \dots$
\choice
{$13\text{ }l$}
{\True $14\text{ }l$}
{$15\text{ }l$}
{$12\text{ }l$}
\loigiai{
$8 + 6 = 14\text{ }l$.\\
\textbf{Đáp án đúng: B}
}
\end{ex}

\begin{ex} % Câu 4
Đường gấp khúc gồm hai đoạn thẳng dài $4\text{ cm}$ và $5\text{ cm}$. Độ dài đường gấp khúc đó là
\choice
{$8\text{ cm}$}
{$10\text{ cm}$}
{$1\text{ cm}$}
{\True $9\text{ cm}$}
\loigiai{
Độ dài đường gấp khúc: $4 + 5 = 9\text{ cm}$.\\
\textbf{Đáp án đúng: D}
}
\end{ex}

\begin{ex} % Câu 5
So sánh: $7 + 5 \dots 6 + 7$.
\choice
{$>$}
{$=$}
{\True $<$}
{Không so sánh được}
\loigiai{
$7 + 5 = 12$; $6 + 7 = 13$. Vì $12 < 13$ nên điền dấu $<$.\\
\textbf{Đáp án đúng: C}
}
\end{ex}

\begin{ex} % Câu 6
Tính: $9 + 8 - 10 = \dots$
\choice
{\True $7$}
{$8$}
{$6$}
{$9$}
\loigiai{
$9 + 8 = 17$; $17 - 10 = 7$.\\
\textbf{Đáp án đúng: A}
}
\end{ex}

\vspace{4mm}
\textbf{\large II. PHẦN TỰ LUẬN (7,0 điểm)}

{\itshape Thí sinh trình bày chi tiết bài làm vào giấy kiểm tra.}

\begin{bt}[1,0 điểm]
Tính nhẩm:
\begin{multicols}{2}
\begin{enumerate}[a)]
  \item $9 + 6$
  \item $7 + 7$
  \item $14 - 8$
  \item $16 - 9$
\end{enumerate}
\end{multicols}
\loigiai{
\begin{itemize}
  \item $9 + 6 = 15$
  \item $7 + 7 = 14$
  \item $14 - 8 = 6$
  \item $16 - 9 = 7$
\end{itemize}
}
\end{bt}

\begin{bt}[1,5 điểm]
Đặt tính rồi tính:
\begin{multicols}{3}
\begin{enumerate}[a)]
  \item $37 + 53$
  \item $89 - 22$
  \item $55 + 42$
\end{enumerate}
\end{multicols}
\loigiai{
\begin{itemize}
  \item $37 + 53 = 90$
  \item $89 - 22 = 67$
  \item $55 + 42 = 97$
\end{itemize}
}
\end{bt}

\begin{bt}[1,0 điểm]
Tính giá trị biểu thức kèm đơn vị đo:
\begin{enumerate}[a)]
  \item $23\text{ kg} + 10\text{ kg} - 10\text{ kg}$
  \item $32\text{ }l - 13\text{ }l + 9\text{ }l$
\end{enumerate}
\loigiai{
\begin{enumerate}[a)]
  \item $23\text{ kg} + 10\text{ kg} - 10\text{ kg} = 33\text{ kg} - 10\text{ kg} = 23\text{ kg}$.
  \item $32\text{ }l - 13\text{ }l + 9\text{ }l = 19\text{ }l + 9\text{ }l = 28\text{ }l$.
\end{enumerate}
}
\end{bt}

\begin{bt}[1,5 điểm]
Bạn Minh có $22$ bông hoa. Bạn Phong có nhiều hơn bạn Minh $5$ bông hoa. Hỏi bạn Phong có bao nhiêu bông hoa?
\loigiai{
Số bông hoa bạn Phong có là:
\[ 22 + 5 = 27 \text{ (bông hoa)} \]
\textbf{Đáp số:} $27$ bông hoa.
}
\end{bt}

\begin{bt}[1,0 điểm]
Một thùng chứa đầy $42\text{ }\text{kg}$ đường kính. Buổi sáng cửa hàng đã bán ra $19\text{ }\text{kg}$. Hỏi trong thùng còn lại bao nhiêu ki-lô-gam đường kính?
\loigiai{
Trong thùng còn lại số ki-lô-gam đường kính là:
\[ 42 - 19 = 23\text{ }\text{kg} \]
\textbf{Đáp số:} $23\text{ }\text{kg}$ đường kính.
}
\end{bt}

\begin{bt}[1,0 điểm]
Tìm số tròn chục lớn nhất có hai chữ số mà tổng các chữ số của nó bằng $8$.
\loigiai{
Số tròn chục có chữ số hàng đơn vị là $0$.\\
Chữ số hàng chục là: $8 - 0 = 8$.\\
Vậy số cần tìm là $80$.\\
\textbf{Đáp số:} $80$.
}
\end{bt}
